\begin{enumerate}
	\item Pokud se zdá, že agent hráče nenásleduje a odchází jedním směrem pryč, nebo se na obrazovce zobrazí chybová hláška znázorněná na obrázku \ref{obrB1}, tato příloha vám poradí, jak danou chybu opravit.
	
\begin{figure}[ht]
	\includegraphics[scale=1.1]{B-1.png} 
	\centering
	\label{obrB1} 
\end{figure}
	
	\item Tato chyba je bohužel známým problémem pluginu NevarokML, ale naštěstí má jednoduché řešení.
	\item Po otevření projektu klikněte ve spodní části editoru na tlačítko \uv{Content Drawer} a přejděte do složky \texttt{All/Content/NPC/AI}, jak je znázorněno na obrázku~\ref{obrB2}.
	\item Otevřete \uv{Blueprint} s názvem \texttt{NPC\_Environment}.
	
\begin{figure}[ht]
	\includegraphics[scale=0.45]{B-2.png} 
	\centering
	\label{obrB2} 
\end{figure}

\newpage

	\item Po otevření \uv{Blueprintu} byste měli na levé straně vidět sekci \uv{Functions}. Zde otevřete funkci \uv{CollectObservations}, jak je znázorněno na obrázku~\ref{obrB3}. Pokud tuto sekci nevidíte, klikněte nahoře na položku \uv{Window} a v části \uv{Graph Editor} zaškrtněte možnost \uv{My Blueprint}.
	
\begin{figure}[ht]
	\includegraphics[scale=1.1]{B-3.png} 
	\centering
	\label{obrB3} 
\end{figure}

	\item Ve funkci \uv{CollectObservations} se nachází uzel \uv{Set Index Counter} a pět uzlů s názvem \uv{Index Plus Plus}. Pokud se agent nepohybuje nebo se zobrazila chyba z obrázku~\ref{obrB1}, je pravděpodobné, že tyto uzly mají nesprávné hodnoty, jak je vidět na obrázku~\ref{obrB4}.

	\item V uzlu \uv{Set Index Counter} nahraďte obě červeně označené hodnoty \texttt{-1} hodnotou \texttt{0}. U všech pěti uzlů \uv{Index Plus Plus} pak změňte hodnoty tak, aby první pole obsahovalo \texttt{0} a druhé \texttt{1}, jak ukazuje obrázek~\ref{obrB5}.
	
	\begin{figure}[ht]
	\includegraphics[scale=1.0]{B-4.png} 
	\centering
	\label{obrB4} 
\end{figure}

\begin{figure}[ht]
	\includegraphics[scale=1.03]{B-5.png} 
	\centering
	\label{obrB5} 
\end{figure}

	\item Tímto je oprava dokončena a agent by nyní měl fungovat správně.

\end{enumerate}